% Simulated Annealing (Heuristics)

A randomized optimization technique used to find global minima/maxima for NP-hard problems (e.g., TSP) by slowly lowering a temperature $T$ parameter to escape local optima.

\begin{lstlisting}
mt19937_64 sa_rng(1337);
inline double get_rand() { return (double)sa_rng() / mt19937_64::max(); }

struct State {
    long long cost;
    void init() { /* Construct initial valid state and calculate cost */ }
    void mutate() { /* Apply a small random perturbation to state */ }
    void revert() { /* Revert the mutation if rejected */ }
};

void simulated_annealing() {
    State current; current.init();
    State best = current;

    double T = 1e5;        // Initial temperature
    double T_min = 1e-3;   // Minimum cut-off temperature
    double alpha = 0.9995; // Cooling schedule factor

    auto start_time = chrono::steady_clock::now();

    while (T > T_min) {
        // Time-limit safety guard (e.g., break at 0.95 seconds to prevent TLE)
        auto current_time = chrono::steady_clock::now();
        if (chrono::duration<double>(current_time - start_time).count() > 0.95) break;

        long long old_cost = current.cost;
        current.mutate();
        long long new_cost = current.cost;

        long long diff = new_cost - old_cost; // Assume minimization problem

        if (diff < 0 || exp(-diff / T) > get_rand()) {
            if (new_cost < best.cost) best = current; // Accept and update best
        } else {
            current.revert(); // Reject mutation
        }
        T *= alpha; // Cool down
    }
}
\end{lstlisting}
